Un radar detecta un avión en $A = (2, 3, 5)$ km y una montaña cuya cara es el plano $\Pi: 3x + 2y + z = 30$ (distancias en km).
\begin{enumerate}
    \item Calcular la distancia perpendicular del avión al plano de la montaña.
    \item El avión vuela en línea recta con velocidad proporcional a $\vec{v} = (1, -1, -2)$. Determinar si su trayectoria intersecta la cara de la montaña; si es así, encontrar el punto de impacto.
    \item Encontrar el plano que contiene la trayectoria del avión y es perpendicular a $\Pi$.
    \item Calcular el ángulo de aproximación: el ángulo agudo entre la dirección de vuelo $\vec{v}$ y el plano $\Pi$.

    \textit{Sugerencia:} El ángulo entre una recta y un plano es el complemento del ángulo entre $\vec{v}$ y $\vec{n}$.
\end{enumerate}